\subsection{The space of Haar measures on the closed subgroups of G}

Lemma 1. --- Let \(\alpha\) be a positive measure \(\neq 0\) on G, S its support; the following two conditions are equivalent:

\begin{enumerate}
\def\labelenumi{\alph{enumi})}
\item
  S is a closed subgroup of G and the measure induced by \(\alpha\) on S is a right Haar measure on S.
\item
  \(\delta(s)\alpha = \alpha\) for every \(s \in S\) .
\end{enumerate}

Moreover, when these conditions are satisfied, the set of \(t \in \mathbf{G}\) such that \(\delta(t)\alpha = \alpha\) is equal to \(S\) .

It is clear that a) implies b); conversely, the relation b) implies that \(Sx = S\) for every \(x \in S\) ; in other words, the relations \(x \in S\) and \(y \in S\) imply that \(y \in Sx\) , or again that \(yx^{-1} \in S\) , and since \(S\) is nonempty, \(S\) is a closed subgroup of \(G\) . The set of \(t \in G\) such that \(St = S\) is then equal to \(S\) itself, whence the last assertion.

For the rest of the section, we denote by \(\Gamma\) the set of positive measures \(\neq0\) on G satisfying the conditions of Lemma 1, and for every \(\alpha\in\Gamma\) we denote by \(H_{\alpha}\) the closed subgroup of G that is the support of \(\alpha\) .

PROPOSITION 1. --- The set \(\Gamma\) is closed in the space \(\mathcal{M}_{+}(G) = \{0\}\) equipped with the vague topology.

We first prove the following lemmas:

Lemma 2. --- Let X be a locally compact space and for every measure \(\alpha \in \mathcal{M}_{+}(X) - \{0\}\) , let \(S_{\alpha}\) be the support of \(\alpha\) . Let \(\Phi\) be a filter on \(\mathcal{M}_{+}(X) - \{0\}\) that converges vaguely to a measure \(\alpha_{0} \neq 0\) . Then, for every neighborhood V of a point s of the support of \(\alpha_{0}\) , there exists a set \(M \in \Phi\) such that, for every \(\alpha \in M\) , one has \(V \cap S_{\alpha} \neq \varnothing\) .

For, if \(\varphi \in \mathcal{K}_{+}(\mathbf{X})\) is a function with support contained in \(\mathbf{V}\) and such that \(\int \varphi(x)d\alpha_0(x) > 0\) , by definition there exists a set \(\mathbf{M} \in \Phi\) such that \(\int \varphi(x)d\alpha(x) > 0\) for all \(\alpha \in \mathbf{M}\) , which implies \(\mathbf{V} \cap \mathbf{S}_{\alpha} \neq \emptyset\) .

Lemma 3. --- Let E be a set filtered by a filter \(\Phi\) , and let \(\xi \mapsto \alpha(\xi)\) be a mapping of E into \(\Gamma\) that converges vaguely with respect to \(\Phi\) to a measure \(\alpha_0 \neq 0\) . On the other hand, let \(\xi \mapsto t_{\xi}\) be a mapping of E into G such that \(t_{\xi} \in \mathrm{H}_{\alpha(\xi)}\) for every \(\xi \in \mathrm{E}\) . If \(s\) is a cluster point of the mapping \(\xi \mapsto t_{\xi}\) with respect to \(\Phi\) , then \(\delta(s)\alpha_0 = \alpha_0\) .

Replacing if necessary \(\Phi\) by a finer filter, we can suppose that \(s\) is a limit of \(\xi \mapsto t_{\xi}\) with respect to \(\Phi\) ; by Lemma 1, \(\delta(t_{\xi})\alpha(\xi) = \alpha(\xi)\) for every \(\xi \in E\) , and the conclusion follows from the continuity of the mapping \((u,\lambda) \mapsto \delta(u)\lambda\) on \(G \times \mathcal{M}_{+}(G)\) (§3, No.~3, Prop. 13).

To prove Prop. 1 it suffices, by Lemma 1, to show that if a filter \(\Psi\) on \(\Gamma\) converges vaguely to a measure \(\alpha_0 \neq 0\) and if \(s\) belongs to the support of \(\alpha_0\) , then \(\delta(s)\alpha_0 = \alpha_0\) . Now, for every neighborhood \(V\) of \(s\) in \(G\) , there exists an \(M \in \Psi\) such that, for every \(\alpha \in M\) , one has \(V \cap H_\alpha \neq \varnothing\) , by Lemma 2. For every neighborhood \(V\) of \(s\) and every \(\alpha \in \Gamma\) , let \(t_{V,\alpha}\) be a point of \(V \cap H_\alpha\) if \(V \cap H_\alpha \neq \varnothing\) , and any point of \(H_\alpha\) in the contrary case; if \(\Theta\) is the section filter of the filter of neighborhoods of \(s\) , and \(\Phi\) is the product filter \(\Theta \times \Psi\) , then \(s\) is, by the foregoing, a cluster point of \((V,\alpha) \mapsto t_{V,\alpha}\) with respect to \(\Phi\) . Since, on the other hand, the mapping \((V,\alpha) \mapsto \alpha\) has \(\alpha_0\) as limit with respect to \(\Phi\) , the proposition follows from Lemma 3.

PROPOSITION 2. --- Let \(\varphi\) be a function in \(\mathcal{K}_{+}(\mathbf{G})\) such that \(\varphi(e) > 0\) . Then the set \(\Gamma_{\varphi}\) of measures \(\alpha \in \Gamma\) such that \(\int \varphi(x) d\alpha(x) = 1\) is compact for the vague topology.

The set \(\Gamma_{\varphi}\) is the intersection of \(\Gamma\) with the hyperplane of \(\mathcal{M}(\mathrm{G})\) formed by the \(\alpha\) such that \(\int \varphi(x)d\alpha(x)=1\) ; since this hyperplane is vaguely closed in \(\mathcal{M}(\mathrm{G})\) and does not contain 0, it follows from Prop. 1 that \(\Gamma_{\varphi}\) is vaguely closed in \(\mathcal{M}(\mathrm{G})\) . It therefore suffices to show that for every compact subset K of G, one has \(\sup_{\alpha\in\Gamma_{\varphi}}\alpha(\mathrm{K})<+\infty\) (Ch. III, §1,

No.~9, Prop. 15). Now, let U be the open neighborhood of e in G defined by the inequality \(\varphi(x) > \varphi(e)/2\) ; since \(1 = \int \varphi(x) d\alpha(x) \geqslant \int_{\mathrm{U}} \varphi(x) d\alpha(x)\) for \(\alpha \in \Gamma_{\varphi}\) , one sees that, on setting \(c = 2/\varphi(e)\) , one has \(\alpha(\mathrm{U}) \leqslant c\) for every \(\alpha \in \Gamma_{\varphi}\) . Let V be a symmetric open neighborhood of e in G such that \(V^{2} \subset U\) ; let us show that \(\alpha(Vx) \leqslant c\) for every \(x \in G\) and every \(\alpha \in \Gamma_{\varphi}\) . Indeed, this relation is trivial if Vx does not intersect the support \(H_{\alpha}\) of \(\alpha\) ; if, on the contrary, there exists an \(h \in Vx \cap H_{\alpha}\) , then h = vx for some \(v \in V\) , whence

\[
\mathbf {V} x = \mathbf {V} v ^ {- 1} h \subset \mathbf {V} ^ {2} h \subset \mathrm{U} h,
\]

and since \(\delta(h)\alpha = \alpha\) , it follows that \(\alpha(\mathrm{V}x) \leqslant \alpha(\mathrm{U}h) = \alpha(\mathrm{U}) \leqslant c\) . Now let \((x_i)_{1 \leqslant i \leqslant n}\) be a sequence of points of K such that the \(Vx_i\) form a covering of K; it follows from the foregoing that \(\alpha(K) \leqslant \sum_{i=1}^{n} \alpha(Vx_i) \leqslant nc\) for every \(\alpha \in \Gamma_\varphi\) ; Q.E.D.

PROPOSITION 3. --- Under the hypotheses of Prop. 2, the mapping \(\alpha \mapsto \left( \langle \varphi, \alpha \rangle, \frac{\alpha}{\langle \varphi, \alpha \rangle} \right)\) is a homeomorphism of \(\Gamma\) onto the product space \(\mathbf{R}_{+}^{*} \times \Gamma_{\varphi}\) .

Since the mapping \(\alpha \mapsto \langle \varphi, \alpha \rangle\) is vaguely continuous, it suffices to observe that \(\langle \varphi, \alpha \rangle \neq 0\) for every measure \(\alpha \in \Gamma\) , since \(e\) belongs to the support \(\mathbf{H}_{\alpha}\) of \(\alpha\) and \(\varphi(e) > 0\) .

